\section{Robustness and Economic Value}\label{sec:robust}

\subsection{Estimators, windows, targets and learners}

Panel A of Table~\ref{tab:robust} varies the construction of the persistence
predictors and the target at the weekly horizon, where model $C$ performs
best. Its gain over HAR remains between 5.2\% and 6.4\% when $d$ is estimated
by local Whittle instead of GPH, over 500 or 1000 days instead of 750, or for
the more or less liquid half of the sample, and it is 2.6\% when the target
is the mean squared return. Each gain is of the same order as that of HAR-X
in the baseline, so none of these choices alters the conclusions of
Section~\ref{sec:results}. A GARCH(1,1) forecast, corrected point-in-time for
its level, is 19.9\% worse than HAR, consistent with the HAR family being the
stronger benchmark for range-based variance.

\begin{table}[H]
\caption{Robustness of model $C$ at $h=5$, and the machine-learning estimators.}
\label{tab:robust}
\footnotesize
\begin{tabularx}{\textwidth}{Xrr}
\toprule
\multicolumn{3}{l}{\textit{Panel A. Model $C$ versus HAR at $h=5$}} \\
\midrule
Variant & $\Delta$MSE vs HAR & DM \\
\midrule
Headline (GPH, $W=750$) & +5.79\% & +4.22$^{***}$ \\
Local Whittle instead of GPH & +5.98\% & +4.13$^{***}$ \\
Window 500 days & +5.33\% & +3.56$^{***}$ \\
Window 1000 days & +6.19\% & +4.16$^{***}$ \\
Target: squared returns & +2.63\% & +2.89$^{***}$ \\
Less liquid half & +6.36\% & +4.16$^{***}$ \\
More liquid half & +5.19\% & +3.98$^{***}$ \\
GARCH(1,1) instead of model $C$ & -19.94\% & -9.04$^{***}$ \\
\midrule
\multicolumn{3}{l}{\textit{Panel B. Model $D$: model $C$'s predictors, other estimators ($\Delta$MSE vs HAR-X at $h=1$ / $5$ / $22$)}} \\
\midrule
Estimator & $\Delta$MSE & DM \\
\midrule
Lasso & -0.8\% / -1.3\% / -3.1\% & -2.6 / -2.2 / -2.6 \\
Ridge & -1.5\% / -1.9\% / -4.2\% & -3.7 / -2.6 / -3.0 \\
Elastic net & -0.8\% / -1.4\% / -3.3\% & -2.5 / -2.2 / -2.5 \\
Random forest & -5.8\% / -6.2\% / -7.0\% & -4.5 / -3.9 / -3.2 \\
Gradient boosting & -6.2\% / -6.1\% / -10.3\% & -6.2 / -5.4 / -4.2 \\
\bottomrule
\end{tabularx}
\noindent{\footnotesize{\textit{Notes:} Panel A: percentage MSE reduction of model $C$ relative to HAR at $h=5$ under each variant; the GARCH row compares a level-corrected GARCH(1,1) forecast with HAR. Panel B: shrinkage hyperparameters and gradient-boosting tree size, leaf size and number of trees chosen by time-series cross-validation inside each training window; the random forest uses 200 trees with a minimum leaf of 20; refit every 20 weeks.}}
\end{table}


Panel B replaces the linear regression on the 18 predictors of model $C$ by
five machine-learning estimators: the lasso, ridge regression and the elastic
net, whose penalties are chosen by time-series cross-validation within each
training window; a random forest; and gradient boosting, whose tree size,
minimum leaf size and number of trees are chosen in the same way. All five
are significantly worse than HAR-X at every horizon. The shrinkage
estimators lose 0.8\% to 4.2\%, and the tree ensembles 5.8\% to 10.3\%. In
this sample, neither regularization nor nonlinearity extracts information
from the persistence predictors that the linear HAR-X model misses. Tuning
matters for this conclusion: gradient boosting with fixed default settings
(400 trees with 31 leaves) loses 12\% to 21\% relative to HAR, an artifact of
overfitting 430 to about 1070 training observations that tuning largely
removes.

\subsection{Volatility-managed portfolios}

A forecast that is statistically no more accurate can still be economically
useful. We examine this with the volatility management of
\citet{moreira2017volatility}, applied stock by stock: each week the
position in stock $i$ is scaled by $c_{i,t}/\hat\sigma^2_{i,t+5\mid t}$,
where $\hat\sigma^2_{i,t+5\mid t}$ is the weekly-horizon variance forecast of
a given model and $c_{i,t}$ is set in real time so that the managed and
unmanaged positions have the same past variance. Positions are funded at the
three-month Treasury bill rate, and returns are measured in excess of it. In
addition to the HAR, HAR-X and model $C$ forecasts, we include a model-free
rule that scales positions by trailing 22-day variance.

\begin{table}[H]
\caption{Volatility-managed portfolios (excess returns, weekly, 2014--2026).}
\label{tab:portfolios}
\footnotesize
\begin{tabularx}{\textwidth}{Xrrrrr}
\toprule
Portfolio & Sharpe & Sharpe, 10 bp & Turnover & Sharpe, COVID & Max drawdown \\
\midrule
Equal weight & 0.76 & 0.76 & 1.1 & 0.88 & -30.6\% \\
Trailing 22-day variance & 0.76 & 0.69 & 6.6 & 1.09 & -15.7\% \\
HAR & 0.75 & 0.64 & 11.5 & 1.29 & -17.3\% \\
HAR-X & 0.76 & 0.66 & 10.5 & 1.38 & -15.6\% \\
Model $C$ & 0.76 & 0.66 & 10.3 & 1.62 & -15.8\% \\
\bottomrule
\end{tabularx}
\noindent{\footnotesize{\textit{Notes:} Each stock's position is scaled by $c_{i,t}/\hat\sigma^2_{i,t+5|t}$ with $c_{i,t}$ set in real time so that the managed and unmanaged positions have equal past variance; positions are funded at the three-month Treasury bill rate; 594 weeks after a 52-week warm-up. Turnover: notional traded per unit of capital per year, including drift. COVID: March to December 2020 (43 weeks). Sharpe differences, Ledoit--Wolf HAC $p$ (block bootstrap $p$): model $C$ vs HAR-X 0.94 (0.93) full sample, 0.10 (0.31) COVID; model $C$ vs equal weight 0.97 (0.98).}}
\end{table}


Table~\ref{tab:portfolios} shows no economic value from any of the forecasts.
Every managed portfolio has a full-sample Sharpe ratio of 0.75 or 0.76, equal
to that of the equal-weighted portfolio, and none of the differences is close
to significant. Management lowers volatility and roughly halves the maximum
drawdown, but it reduces expected excess returns in proportion. The managed
portfolios trade about ten times their capital each year, and at a cost of
10 basis points per unit traded their Sharpe ratios fall to 0.64--0.66, below
that of equal weighting. The one period in which model $C$ stands out, March
to December 2020, contains 43 weekly observations, and its advantage over
HAR-X in that period (a Sharpe ratio of 1.62 against 1.38) is not
significant. These results are in line with the evidence that volatility
management of factor portfolios is fragile out of sample
\citep{cederburg2020volatility, barroso2021volatility}.
